\documentclass[11pt,a4paper]{article}

\usepackage[utf8]{inputenc}
\usepackage[T1]{fontenc}
\usepackage{lmodern}
\usepackage[margin=2.5cm]{geometry}
\usepackage{amsmath,amssymb}
\usepackage{graphicx}
\usepackage{booktabs}
\usepackage{array}
\usepackage{xcolor}
\usepackage{hyperref}
\usepackage{caption}
\usepackage{subcaption}
\usepackage{enumitem}

\hypersetup{
  colorlinks=true,
  linkcolor=blue!60!black,
  citecolor=blue!60!black,
  urlcolor=blue!60!black
}

\graphicspath{{figures/}{../outputs/tpv/}{../outputs/unet/}{../outputs/diffpir/}{../outputs/comparison/}}

\newcommand{\placeholder}[1]{%
  \fbox{%
    \begin{minipage}[c][0.22\textheight][c]{0.88\linewidth}
      \centering
      \textbf{Figure placeholder}\\[0.5em]
      \small \texttt{\detokenize{#1}}
    \end{minipage}%
  }%
}

\newcommand{\maybeincludegraphics}[2][]{%
  \IfFileExists{#2}{\includegraphics[#1]{#2}}{\placeholder{#2}}%
}

\newcommand{\todo}[1]{\textcolor{red!70!black}{\textbf{TODO:} #1}}

\title{\textbf{Sparse-View CT Reconstruction on the Mayo Dataset}\\
\large Variational, End-to-End, and Generative Approaches}
\author{Group B -- Computational Imaging Project}
\date{\today}

\begin{document}
\maketitle

\begin{abstract}
This report studies sparse-view computed tomography (CT) reconstruction as a
computational imaging inverse problem. The same degraded Mayo CT measurements
are used throughout: images resized to $256\times256$, parallel-beam projection
geometries with $180$, $90$, $60$, and $45$ views, and additive Gaussian
measurement noise with relative level $0.005$. Three methodological families are
compared: Total $p$-Variation (TpV) regularization---a classical model-based
approach that encodes an explicit sparsity prior on image gradients; a
Generalized ResUNet post-processing network---a supervised end-to-end method
trained from filtered back-projection (FBP) proxy images; and a DiffPIR-based
diffusion reconstruction method---a plug-and-play generative prior that guides a
learned reverse diffusion process with CT data-consistency updates. Quantitative
evaluation is performed through PSNR and SSIM, complemented by visual comparison
panels across the four sparse-view configurations.
\end{abstract}

% ============================================================
\section{Introduction}
% ============================================================

Computational imaging is the joint design of image acquisition, the physical
forward model, and computational reconstruction algorithms in order to recover
information that a direct observation cannot provide
\cite{mccollough2016lowdose}. Computed tomography is one of its most prominent
instances: the internal attenuation map of an object is estimated from a
collection of indirect X-ray line-integral measurements.

\paragraph{Forward model.}
In the two-dimensional setting of this project, the unknown image is
$x\in\mathbb{R}^{256\times256}$ and the acquisition follows the discrete
linear forward model
\begin{equation}
  y^{\delta} = K_\theta x + e,
  \label{eq:ct-forward-model}
\end{equation}
where $K_\theta:\mathbb{R}^{N}\to\mathbb{R}^{M}$ is the parallel-beam
Radon-type projection operator parameterized by the set of projection angles
$\theta$, $y^{\delta}\in\mathbb{R}^{M}$ is the noisy measured sinogram, and
$e\in\mathbb{R}^{M}$ is additive measurement noise. The operator $K_\theta$
is a large, severely ill-conditioned matrix: its singular values decay rapidly
to zero, so small perturbations in the data $y^{\delta}$ produce arbitrarily
large errors in any naive inversion. This property is the mathematical signature
of an \emph{ill-posed} inverse problem in the sense of Hadamard---existence,
uniqueness, or stability of the solution may fail.

\paragraph{Sparse-view task.}
The project addresses the sparse-view CT scenario, in which the number of
projection angles is strongly reduced relative to the Nyquist-Shannon angular
sampling rate required for exact inversion:
\begin{equation}
  n_\theta \in \{180,\,90,\,60,\,45\}.
  \label{eq:view-settings}
\end{equation}
As $n_\theta$ decreases, the operator $K_\theta$ loses rows, reducing the
measured data and increasing the degree of under-determination. In practice, a
direct filtered back-projection (FBP) reconstruction from sparse-view data
exhibits characteristic streak artifacts that grow more severe as $n_\theta$
decreases. Regularization---either explicit through a variational penalty or
implicit through a learned prior---is therefore indispensable.

\paragraph{Methodological families.}
Three reconstruction families are studied, each corresponding to one of the
principal strategies introduced in the course:
\begin{enumerate}[leftmargin=*]
  \item \textbf{Model-based variational approach}: Total $p$-Variation (TpV)
        regularization \cite{rudin1992nonlinear}, which expresses prior
        knowledge as an explicit penalty on the image gradient magnitude.
  \item \textbf{End-to-end supervised learning}: a Generalized ResUNet
        post-processor \cite{ronneberger2015unet} that maps an FBP proxy
        directly to a clean reconstruction by minimizing a supervised loss over
        training data.
  \item \textbf{Diffusion-based generative prior}: DiffPIR
        \cite{zhu2023denoising}, which embeds the CT measurement constraint
        inside the reverse denoising trajectory of a diffusion model.
\end{enumerate}

\subsection*{Mayo Dataset}

The dataset used is the Mayo CT dataset provided in the course material
\cite{mccollough2016lowdose}. PNG slices are organized by patient folders and
split at patient level to avoid data leakage between train and test sets.
The processed split used by the project notebooks is summarized in
Table~\ref{tab:mayo-split}. Each slice is resized to $256\times256$ using
bilinear interpolation and normalized to $[0,1]$ before the CT forward model
is applied. Processed data are saved as one \texttt{.pt} file per patient,
preserving the correspondence between patient identifiers and slices.

\begin{table}[t]
  \centering
  \caption{Mayo processed dataset organization used by the project notebooks.}
  \label{tab:mayo-split}
  \begin{tabular}{lrr}
    \toprule
    Split & Patients & Slices \\
    \midrule
    Train & 8 & 2585 \\
    Validation & 2 & 721 \\
    Test & 1 & 327 \\
    \bottomrule
  \end{tabular}
\end{table}

For each clean image $x$, four degraded sinograms are generated, one per view
configuration. The noise model is applied in the measurement domain: if
$y = K_\theta x$ is the clean sinogram, the noisy data are constructed as
\begin{equation}
  y^{\delta} = y + e,
  \qquad
  \frac{\|e\|_2}{\|y\|_2} \approx 0.005,
  \label{eq:relative-noise}
\end{equation}
where $e$ is sampled from a Gaussian distribution and scaled to match the
prescribed relative level. Fixing the same noise realization for all methods
guarantees a fair comparison: different noise instances would make numerical
results non-comparable.

\subsection*{Tools}

The implementation uses the IPPy library distributed with the course material.
The main components employed are:
\begin{itemize}[leftmargin=*]
  \item \path{IPPy.operators.CTProjector}: wraps the ASTRA-based parallel-beam
        forward operator $K_\theta$ and its adjoint $K_\theta^\top$
        \cite{astra2016};
  \item \path{IPPy.solvers.FBP}: computes the Ram-Lak filtered back-projection
        proxy $x_{\mathrm{FBP}}$, used as input to the ResUNet and as a
        baseline for comparison;
  \item \path{IPPy.solvers.ChambollePockTpVUnconstrained}: implements the
        primal-dual TpV solver via Chambolle--Pock iterations
        \cite{chambolle2011first};
  \item \path{IPPy.nn.models.UNet}: provides the ResUNet architecture used for
        end-to-end post-processing \cite{ronneberger2015unet};
  \item \path{IPPy.utilities.metrics}: computes PSNR and SSIM \cite{wang2004ssim}.
\end{itemize}

\begin{figure}[t]
  \centering
  \maybeincludegraphics[width=0.95\linewidth]{dataset_sanity_check.png}
  \caption{One normalized Mayo CT slice and the corresponding noisy
    sparse-view sinograms for $180$, $90$, $60$, and $45$ projection angles.}
  \label{fig:dataset-sanity-check}
\end{figure}

% ============================================================
\section{Selected Methods}
% ============================================================

\subsection{Total \texorpdfstring{$p$}{p}-Variation Regularization}

\paragraph{Variational formulation.}
The first selected method is a model-based variational approach. Given the
ill-posed nature of the CT inverse problem, a principled remedy is to replace
the naive inversion by a regularized variational problem that balances data
fidelity with a prior on the image. The general \emph{Tikhonov-type}
formulation is
\begin{equation}
  \widehat{x}
  =
  \arg\min_{x \ge 0}
  \frac{1}{2}\|K_\theta x - y^{\delta}\|_2^2
  + \lambda\, R(x),
  \label{eq:variational-general}
\end{equation}
where the first term enforces consistency with the measured sinogram and the
second introduces prior information on the solution. The regularization
parameter $\lambda > 0$ controls the trade-off: small $\lambda$ favors
data-fidelity and may leave noise and streak artifacts; large $\lambda$
enforces strong regularity but risks over-smoothing and loss of fine anatomical
detail.

\paragraph{Total $p$-Variation prior.}
In this project, the regularizer is the Total $p$-Variation functional
\cite{rudin1992nonlinear}:
\begin{equation}
  R_p(x)
  =
  \sum_i \|\nabla x_i\|_2^p,
  \qquad
  0.1 < p < 0.5.
  \label{eq:tpv-regularizer}
\end{equation}
Here $\nabla x_i \in \mathbb{R}^2$ is the discrete image gradient at pixel $i$,
computed via finite differences. The standard total variation (TV) corresponds
to $p=1$, which yields a convex functional and is well known to promote
piecewise-constant images while preserving sharp edges. By setting $p < 1$,
the TpV penalty is \emph{non-convex} and promotes \emph{sparser} image
gradients than TV: only a few pixels are allowed to have large gradient
magnitudes, while the majority are forced to zero. This property aligns well
with the structure of CT anatomical images, which typically exhibit smooth
regions separated by clinically meaningful edges (organ boundaries,
bone surfaces). The project specifications require $p \in (0.1, 0.5)$; we
use $p = 0.35$.

\paragraph{Chambolle--Pock primal-dual solver.}
Because $R_p$ is non-convex for $p < 1$, the optimization problem
\eqref{eq:variational-general}--\eqref{eq:tpv-regularizer} cannot be solved
by standard convex proximal algorithms. IPPy handles this through a
\emph{majorization-minimization} (reweighted) strategy: the TpV term is
linearized around the current estimate by introducing pixel-wise weights
\begin{equation}
  W^{(k)}(x)
  =
  \left(
  \frac{\sqrt{\eta^2 + |\nabla x^{(k)}|^2}}{\eta}
  \right)^{p-1},
  \label{eq:tpv-weight}
\end{equation}
where $\eta > 0$ is a small smoothing parameter that avoids singularities at
zero-gradient locations. At each outer iteration $k$, the weights $W^{(k)}$ are
fixed and a weighted TV subproblem is solved via the
\emph{Chambolle--Pock primal-dual} scheme \cite{chambolle2011first}. The
algorithm alternates between a gradient-ascent step in the dual variable
(associated with the TV-like term) and a gradient-descent step in the primal
variable (associated with data fidelity), with step sizes governed by the
operator norms of $K_\theta$ and $\nabla$. Since $p - 1 < 0$, pixels with
larger gradient magnitudes receive smaller weights in subsequent iterations,
which progressively concentrates the sparsity pattern. The reconstruction is
constrained to non-negative values at every step.

\subsection{Generalized ResUNet Post-Processing}

\paragraph{FBP proxy and learned correction.}
The second selected method is a supervised end-to-end approach. Rather than
operating on the sinogram directly, the method first computes a filtered
back-projection reconstruction from the sparse-view noisy data, then trains a
neural network to map the FBP image to the clean slice. The two-stage pipeline
is
\begin{equation}
  x_{\mathrm{FBP}} = \operatorname{FBP}(y^{\delta}),
  \qquad
  \widehat{x} = f_\Theta(x_{\mathrm{FBP}}),
  \label{eq:resunet-mapping}
\end{equation}
where $f_\Theta$ is the learned ResUNet and $\Theta$ denotes its trainable
parameters. The FBP step is not a standalone reconstruction method in this
context: it is used only to convert the sinogram to the image domain, making
the learning problem an image-to-image regression task rather than a
sinogram-to-image problem.

\paragraph{Supervised MSE loss.}
The network is trained by minimizing the mean squared error between prediction
and the corresponding clean Mayo slice:
\begin{equation}
  \mathcal{L}(\Theta)
  =
  \frac{1}{N}
  \sum_{j=1}^{N}
  \left\|
  f_\Theta(x_{\mathrm{FBP},j}) - x_j
  \right\|_2^2.
  \label{eq:mse-loss}
\end{equation}
Minimizing the $\ell^2$ loss is equivalent, under a Gaussian likelihood model,
to computing the maximum-likelihood estimate of the network output. In imaging,
MSE training is known to produce outputs that are good in expectation but may
appear slightly over-smoothed, because the network averages over residual
uncertainty in the target distribution.

\paragraph{ResUNet architecture.}
The network architecture is the U-Net \cite{ronneberger2015unet}, extended with
residual blocks in the encoder and decoder branches. The encoder extracts
increasingly abstract features through successive downsampling blocks, while the
decoder reconstructs a full-resolution output using transposed convolutions. Skip
connections between encoder and decoder levels concatenate feature maps of the
same spatial resolution, preserving fine spatial detail that would otherwise be
lost during downsampling. The residual connections inside each block ease gradient
flow during backpropagation and allow the network to learn additive corrections
to an identity mapping, which is particularly appropriate here: the FBP input is
already a reasonable (if artifact-corrupted) image, and the network only needs to
learn a \emph{residual correction} rather than a full synthesis.

\paragraph{Generalized multi-view training.}
A single ResUNet is trained across all four sparse-view configurations. During
each training epoch, the dataloader cycles over all patients and, for each
patient, iterates over all four angle keys $\{180, 90, 60, 45\}$. Each clean
slice is therefore paired with four different FBP inputs corresponding to
different levels of angular undersampling. This strategy produces a
\emph{generalized} artifact-reduction model that is exposed to the full spectrum
of undersampling artifacts during training and can, in principle, perform well
across all view counts without retraining.

\subsection{DiffPIR-Based Diffusion Reconstruction}

\paragraph{Diffusion models as image priors.}
The third selected method is a generative approach based on denoising diffusion
probabilistic models (DDPMs). A diffusion model defines a probability
distribution over clean images through two coupled Markov chains. The
\emph{forward} (noising) process gradually corrupts a clean image $x_0$
towards pure Gaussian noise:
\begin{equation}
  x_t
  =
  \sqrt{\bar{\alpha}_t}\,x_0
  +
  \sqrt{1-\bar{\alpha}_t}\,\epsilon,
  \qquad
  \epsilon \sim \mathcal{N}(0,I),
  \label{eq:diffusion-forward}
\end{equation}
where $t \in \{1,\ldots,T\}$ indexes the diffusion step and
$\bar{\alpha}_t = \prod_{s=1}^{t}(1-\beta_s)$ is the cumulative product of the
noise schedule $\{\beta_s\}$. At $t=T$, $\bar{\alpha}_T \approx 0$ and $x_T$
is approximately standard Gaussian noise. A neural network $\epsilon_\varphi$
is trained to predict the noise component $\epsilon$ from the noisy image $x_t$
and the time step $t$; equivalently, it predicts the \emph{score}
$\nabla_{x_t}\log p(x_t)$ of the data distribution.

\paragraph{DiffPIR: diffusion for plug-and-play inversion.}
DiffPIR \cite{zhu2023denoising} adapts the reverse diffusion sampling loop for
solving inverse problems. At each reverse step $t$, after the standard
denoising update, a \emph{data-consistency projection} is applied to enforce
agreement with the measured CT sinogram. The data-fidelity term is
\begin{equation}
  \mathcal{D}(x)
  =
  \frac{1}{2}\|K_\theta x - y^{\delta}\|_2^2,
  \label{eq:diffpir-data-consistency}
\end{equation}
and its gradient with respect to $x$ is
$\nabla_x \mathcal{D}(x) = K_\theta^\top(K_\theta x - y^{\delta})$.
The data-consistency update is a gradient step of the form
\begin{equation}
  x \leftarrow x - \zeta\, K_\theta^\top(K_\theta x - y^{\delta}),
  \label{eq:diffpir-gradient-step}
\end{equation}
where $\zeta > 0$ is a step-size parameter that controls the strength of the
data correction. This step is interleaved with the learned reverse-diffusion
denoising step at each time $t$. Intuitively, the diffusion prior explores the
manifold of natural-looking CT images, while the data-consistency step
continuously steers the trajectory toward images that are compatible with the
measured sparse-view sinogram.

\paragraph{Trade-off between prior and data.}
The parameter $\zeta$ governs a fundamental trade-off. If $\zeta$ is too small,
the generative prior dominates: the output may be anatomically plausible and
visually sharp, but it may hallucinate structures not present in the measured
data (a critical failure mode in medical imaging). If $\zeta$ is too large, the
data-consistency correction overwhelms the prior and the reconstruction may
reintroduce the original sparse-view streak artifacts, effectively reverting to
an FBP-like result. The number of reverse diffusion steps $T$ and the noise
schedule are additional degrees of freedom: fewer steps accelerate inference
but may degrade image quality; the schedule determines how fast the model
transitions from the learned distribution to the noisy initial state.

% ============================================================
\section{Implementation}
% ============================================================

\subsection{Implementation of Total \texorpdfstring{$p$}{p}-Variation}

The TpV pipeline is implemented in
\texttt{notebooks/01\_TpV\_reconstruction.ipynb}. The notebook loads the
processed Mayo patient file for the test split and selects the fixed central
slice across all view settings for consistency. For each view configuration, a
CT projector is instantiated as
\begin{equation}
  \begin{aligned}
  K_\theta
  &=
  \operatorname{CTProjector}
  \bigl(
  \texttt{img\_shape}=(256,256),\\
  &\hspace{2.2cm}
  \texttt{angles}=\operatorname{linspace}(0,\pi,n_\theta),\;
  \texttt{det\_size}=256,\;
  \texttt{geometry}=\texttt{parallel}
  \bigr).
  \end{aligned}
  \label{eq:ippy-projector}
\end{equation}
The corresponding TpV solver is then created through
\path{IPPy.solvers.ChambollePockTpVUnconstrained}. The
final hyperparameters selected for the test evaluation are
\begin{equation}
  \lambda = 0.01,
  \qquad
  p = 0.35,
  \qquad
  \texttt{maxiter}=150,
  \qquad
  \texttt{tolf}=\texttt{tolx}=10^{-5}.
  \label{eq:tpv-parameters}
\end{equation}
The solver is initialized from a zero image (\texttt{starting\_point=None}).
The ground-truth clean slice is provided only for diagnostic computation of
PSNR and SSIM at each iteration; it does not influence the optimization. The
non-negativity constraint $x \ge 0$ is enforced internally by the solver.

The notebook saves one visual panel per view configuration: each panel contains
the ground-truth image, the TpV reconstruction (normalized to $[0,1]$), and the
absolute error map $|x_{\mathrm{TpV}} - x|$. Four convergence plots are also
saved, containing the objective value, the primal-dual residual, the PSNR
curve, and the SSIM curve as functions of iteration number.

\subsection{Implementation of Generalized ResUNet}

The learned method is implemented in
\texttt{notebooks/02\_ResUnet\_reconstruction.ipynb}. The FBP proxy is
computed in memory for each sinogram batch using
\path{IPPy.solvers.FBP} applied to the saved noisy sinograms. The resulting
FBP images are normalized per image and treated as the degraded input to the
network; they are not used as a standalone reconstruction method.

The network is instantiated through \path{IPPy.nn.models.UNet} with the
following configuration:
\begin{equation}
  \begin{aligned}
  \texttt{ch\_in} &= 1,\quad \texttt{ch\_out} = 1,\\
  \texttt{middle\_ch} &= (32,\,64,\,128,\,256,\,512),\\
  \texttt{n\_layers\_per\_block} &= 2,\\
  \texttt{final\_activation} &= \texttt{None}.
  \end{aligned}
  \label{eq:resunet-config}
\end{equation}
The final layer produces unconstrained logits; outputs are clamped to $[0,1]$
only at evaluation time and for display. This choice avoids the vanishing
gradient issue associated with a sigmoid activation during training.

Training is supervised with the MSE loss. For each patient batch, the
dataloader cycles over all four angle configurations $\{180,90,60,45\}$,
pairing each sinogram with the corresponding clean slice. The training
hyperparameters are
\begin{equation}
  \texttt{learning\_rate}=10^{-4},
  \qquad
  \texttt{batch\_size}=8,
  \qquad
  \texttt{epochs}=2,
  \qquad
  \texttt{patience}=3.
  \label{eq:resunet-training-parameters}
\end{equation}
The optimizer is Adam. The checkpoint
\texttt{weights/unet/resunet\_generalized\_latest.pt} stores the model state,
optimizer state, loss history, validation loss history, and configuration.

\subsection{Implementation of DiffPIR}

The repository configuration selects DiffPIR as the required generative method
for the two-student group. No complete DiffPIR reconstruction notebook is
currently present in the repository; this subsection defines the theoretical
slot that must be completed once the implementation is finalized.

The final DiffPIR implementation must use the same processed Mayo tensors and
the same IPPy CT projector as the other methods. For each view count, the
algorithm should:
\begin{enumerate}[leftmargin=*]
  \item Load the fixed noisy sinogram $y^{\delta}$ for the shared central test slice.
  \item Instantiate \texttt{IPPy.operators.CTProjector} to provide $K_\theta$
        and $K_\theta^\top$.
  \item Initialize the reverse diffusion at $x_T \sim \mathcal{N}(0,I)$.
  \item At each reverse step $t = T, T-1, \ldots, 1$:
        (a) apply the learned denoising step using $\epsilon_\varphi$;
        (b) apply the data-consistency gradient step
            \eqref{eq:diffpir-gradient-step} with weight $\zeta$.
  \item Output $x_0$ as the final reconstruction.
\end{enumerate}
The parameters that must be reported are: the pre-trained denoiser
$\epsilon_\varphi$ used, the total number of reverse diffusion steps $T$, the
noise schedule $\{\beta_t\}$, and the data-consistency weight $\zeta$.

\begin{figure}[t]
  \centering
  \maybeincludegraphics[width=0.9\linewidth]{method_pipeline.png}
  \caption{Complete experimental pipeline: from a Mayo CT slice to the
    sparse-view sinogram, then to TpV, ResUNet, and DiffPIR reconstructions.}
  \label{fig:method-pipeline}
\end{figure}

% ============================================================
\section{Results and Discussion}
% ============================================================

\subsection{Hyperparameter Tuning}

\paragraph{TpV.}
Three parameters govern the TpV reconstruction: the regularization weight
$\lambda$, the gradient-sparsity exponent $p$, and the maximum number of
Chambolle--Pock iterations. The parameter $\lambda$ encodes the relative
confidence in the prior versus the data. A value of $\lambda$ that is too small
may leave streak artifacts and measurement noise; a value that is too large may
over-smooth the image, erasing clinically relevant edges and introducing
piecewise-constant staircasing artifacts (a well-known artifact of TV-type
regularization). The exponent $p$ controls non-convexity: values closer to $0$
enforce more aggressive gradient sparsity but make the landscape increasingly
non-convex and the optimization less reliable; values closer to $1$ approach
standard convex TV. The selected value $p=0.35$ lies inside the required range
$(0.1, 0.5)$ and provides a balance between sparsity and numerical stability.
The number of iterations is set to $150$, with early-stopping tolerance
$\texttt{tolf}=\texttt{tolx}=10^{-5}$.

\paragraph{ResUNet.}
The key hyperparameters are the learning rate, batch size, and architecture
configuration. A learning rate of $10^{-4}$ is used with the Adam optimizer;
the batch size is 8. The architecture uses $5$ encoder-decoder levels with base
channel count $32$ (doubling at each level, reaching $512$ at the bottleneck),
and $2$ residual layers per block. Training proceeds for up to $2$ epochs with
a patience of $3$ epochs for early stopping on validation MSE. These values
reflect a preliminary configuration that can be extended as computational
resources allow.

\paragraph{DiffPIR.}
The tuning process must be completed after implementation. The key parameter is
the data-consistency weight $\zeta$ in the gradient step
\eqref{eq:diffpir-gradient-step}. Additional parameters include the number of
reverse diffusion steps $T$ (typically $100$--$1000$) and the noise schedule.
The expected trade-off is between generative image quality (controlled by the
denoising steps) and physical consistency with the measured CT sinogram
(controlled by $\zeta$).

\subsection{Evaluation Metrics}

Quantitative evaluation uses PSNR and SSIM, computed on images normalized to
$[0,1]$. The PSNR is defined as
\begin{equation}
  \operatorname{PSNR}(\widehat{x},x)
  =
  -20\log_{10}
  \!\left(
  \sqrt{\operatorname{MSE}(\widehat{x},x)}
  \right),
  \qquad
  \operatorname{MSE}(\widehat{x},x)
  =
  \frac{1}{N}
  \sum_{i=1}^{N}
  (\widehat{x}_i - x_i)^2,
  \label{eq:psnr}
\end{equation}
where the images are normalized so that the maximum intensity is $1$. Higher
PSNR (measured in decibels, dB) indicates lower mean squared error. The
formula is equivalent to $10\log_{10}(1/\operatorname{MSE})$ for unit-range
images.

The SSIM \cite{wang2004ssim} is a perceptual metric that combines luminance,
contrast, and structural similarity:
\begin{equation}
  \operatorname{SSIM}(\widehat{x},x)
  =
  \frac{(2\mu_{\widehat{x}}\mu_x + c_1)(2\sigma_{\widehat{x}x} + c_2)}
       {(\mu_{\widehat{x}}^2 + \mu_x^2 + c_1)(\sigma_{\widehat{x}}^2 + \sigma_x^2 + c_2)},
  \label{eq:ssim}
\end{equation}
where $\mu$, $\sigma^2$ denote local means and variances, $\sigma_{\widehat{x}x}$
is the local cross-covariance, and $c_1, c_2$ are small stabilization
constants. SSIM takes values in $[-1,1]$, with $1$ indicating a perfect match.
Unlike PSNR, SSIM is sensitive to local structural patterns and is therefore a
better predictor of perceived reconstruction quality in medical imaging.

Both metrics are computed on the shared central test slice to enable direct
comparison across methods and view counts.

\subsection{Quantitative Results}

Table~\ref{tab:quantitative-results} collects the PSNR and SSIM values for all
methods and view configurations. The ResUNet entries reflect the preliminary
scores printed by \texttt{notebooks/02\_ResUnet\_reconstruction.ipynb} on the
central test slice after only $2$ training epochs. The identical values across
all view counts (PSNR $13.84$ dB, SSIM $0.1259$) indicate that the current
checkpoint is undertrained and has not yet converged to a configuration that
exploits view-specific artifact information; additional epochs are expected to
improve performance and differentiate results across view counts. TpV values
must be filled from \texttt{notebooks/01\_TpV\_reconstruction.ipynb}, and
DiffPIR values after the third method is implemented.

\begin{table}[t]
  \centering
  \caption{PSNR and SSIM comparison on the shared Mayo test slice.}
  \label{tab:quantitative-results}
  \begin{tabular}{llcc}
    \toprule
    Views & Method & PSNR (dB) & SSIM \\
    \midrule
    180 & TpV & \todo{fill} & \todo{fill} \\
    180 & Generalized ResUNet & 13.84 & 0.1259 \\
    180 & DiffPIR & \todo{fill after implementation} & \todo{fill} \\
    90 & TpV & \todo{fill} & \todo{fill} \\
    90 & Generalized ResUNet & 13.84 & 0.1259 \\
    90 & DiffPIR & \todo{fill after implementation} & \todo{fill} \\
    60 & TpV & \todo{fill} & \todo{fill} \\
    60 & Generalized ResUNet & 13.84 & 0.1259 \\
    60 & DiffPIR & \todo{fill after implementation} & \todo{fill} \\
    45 & TpV & \todo{fill} & \todo{fill} \\
    45 & Generalized ResUNet & 13.84 & 0.1259 \\
    45 & DiffPIR & \todo{fill after implementation} & \todo{fill} \\
    \bottomrule
  \end{tabular}
\end{table}

\subsection{Visual Results and Critical Discussion}

The visual comparison is organized as follows. For each method and each view
count, a panel shows (i) the ground-truth slice, (ii) the reconstruction, and
(iii) the absolute error map $|\widehat{x} - x|$. Convergence curves for TpV
(objective, residual, PSNR, SSIM vs.\ iteration) are shown separately. A
global comparison panel summarizes PSNR and SSIM across methods and view counts.

\begin{figure}[t]
  \centering
  \begin{subfigure}{0.48\linewidth}
    \centering
    \maybeincludegraphics[width=\linewidth]{comparison_panel_180.png}
    \caption{180 views}
  \end{subfigure}
  \hfill
  \begin{subfigure}{0.48\linewidth}
    \centering
    \maybeincludegraphics[width=\linewidth]{comparison_panel_45.png}
    \caption{45 views}
  \end{subfigure}
  \caption{Representative comparison panels for high-view and low-view sparse
    CT settings. Each panel shows the ground truth, TpV reconstruction, ResUNet
    reconstruction, and their respective absolute error maps.}
  \label{fig:comparison-panels}
\end{figure}

\begin{figure}[t]
  \centering
  \maybeincludegraphics[width=0.85\linewidth]{project_methods_quantitative_comparison_all_views.png}
  \caption{Aggregate PSNR and SSIM across view counts and methods.}
  \label{fig:quantitative-plot}
\end{figure}

The final critical discussion must address the following points once all
results are available.

\paragraph{Effect of view reduction.}
It should be verified that reconstruction quality (PSNR and SSIM) decreases
monotonically as $n_\theta$ decreases from $180$ to $45$. The FBP proxy, which
is used as the ResUNet input and provides a direct measure of measurement
quality, is expected to degrade significantly at $45$ views.

\paragraph{TpV analysis.}
TpV should be evaluated in terms of: (i) edge preservation, since the
non-convex gradient penalty favors sharp boundaries; (ii) noise suppression,
since the data-fidelity term limits amplification of measurement noise;
(iii) staircasing artifacts, a known pathology of TV-type regularization where
smooth intensity ramps are approximated by staircase functions; and (iv)
sensitivity to $\lambda$ and $p$. In particular, the improvement of TpV over
FBP is expected to grow as $n_\theta$ decreases, because the regularization
prior becomes more beneficial when the data alone provide less information.

\paragraph{ResUNet analysis.}
The ResUNet should be discussed in terms of: (i) artifact removal, assessing
whether the network successfully eliminates streak artifacts; (ii) possible
over-smoothing, a known consequence of MSE-loss training; (iii) generalization
across view counts, verifying whether a single model trained on all four
configurations achieves competitive performance at each; and (iv) dependence on
the FBP proxy quality, since the input to the network degrades as $n_\theta$
decreases. The identical PSNR and SSIM values across all view settings in the
current checkpoint (Table~\ref{tab:quantitative-results}) suggest that the
model has not yet learned to differentiate artifact patterns and requires
additional training epochs to converge.

\paragraph{DiffPIR analysis.}
Once implemented, DiffPIR should be evaluated by assessing whether the
generative prior improves visual quality relative to TpV and ResUNet without
violating CT data consistency. A qualitative check should verify that the
reconstruction does not hallucinate anatomical structures that are absent from
the measured sinogram. The DiffPIR result is expected to be particularly
beneficial at the lowest view count ($45$ views), where the FBP proxy is most
severely corrupted and the generative prior has the most opportunity to fill in
missing information.

% ============================================================
\section{Conclusions}
% ============================================================

This report formulates sparse-view CT reconstruction on the Mayo dataset as an
ill-posed linear inverse problem and presents a structured comparison of three
methodological families. Each family represents a distinct philosophy for
incorporating prior knowledge about the unknown image.

The TpV method provides an interpretable model-based strategy: the prior is an
explicit non-convex functional on the image gradient, encoding the assumption
that CT anatomical images have sparse gradient support. The Chambolle--Pock
primal-dual solver handles the resulting non-convex optimization via
majorization-minimization with adaptive weights. This approach requires no
training data and its behavior is fully determined by the choice of $\lambda$
and $p$; however, it requires per-image iterative optimization and its
expressiveness is limited to the gradient-sparsity assumption.

The Generalized ResUNet provides a supervised end-to-end strategy: the prior is
implicit and distributed across the millions of trainable parameters of the
network. By training on pairs of FBP inputs and clean targets across all four
view configurations, the network learns a data-driven correction function that
is fast to evaluate at test time. Its limitation is dependence on the
distribution of training data: degradations unseen during training may not be
corrected effectively.

The DiffPIR method embeds the CT data-consistency constraint into the reverse
trajectory of a score-based diffusion model, combining a powerful generative
prior with explicit physics-based guidance. Its implementation and quantitative
results remain to be completed.

The final version of the report should insert the missing PSNR and SSIM values,
add the generated figures, and replace the DiffPIR placeholders with exact
implementation details and experimental results. Once complete, the discussion
will compare the three approaches in terms of numerical accuracy, visual
artifact characteristics, computational cost, and robustness across sparse-view
configurations.

\bibliographystyle{unsrt}
\bibliography{references}

\end{document}
